\chapter{Introduction}

\section{The question}

The question worth asking about parallel programming models is not whether
OpenMP is fast. It is what each model costs on identical work. That question is
hard to answer honestly because the usual comparison changes two things at once:
a study that reimplements an algorithm per model ends up measuring the care that
went into each reimplementation, and one that compares different algorithms on
different devices ends up measuring the algorithms.

This project answers it by construction. One numerical library, twelve iterative
solvers written exactly once against a single problem interface, and seven
execution backends behind a single parallelism interface. A solver does not know
which backend it is running on. A backend does not know which solver it is
running. What varies between measurements is therefore the execution model, and
that is enforced by a test rather than promised in a paragraph.

\section{The invariant that makes it work}

The design commits to something stronger than the usual claim of numerical
agreement. With the deterministic reduction mode:

\begin{itemize}
  \item every shared memory backend, at every worker count, produces
        \emph{bit identical} iterates, not merely close ones;
  \item the GPU sweeps produce bit identical values to the CPU as well;
  \item across MPI rank counts, results agree to reduction tolerance, and the
        reason they cannot be bit identical is documented rather than hidden.
\end{itemize}

This is achievable because floating point addition is not associative and the
library therefore refuses to let the association vary. The chunk grid used by
every reduction is fixed by the problem size alone and never by the worker
count, so the partial sums are always combined in the same order regardless of
which worker produced which one or when. Chapter \ref{ch:parallelisation}
explains the mechanism; the equivalence suite asserts exact equality across
twelve solvers, two problem families, four backends and seven worker counts.

Getting there required disabling fused multiply add contraction on both the host
and the device. The SOR update has the shape $ab + cd$, which a compiler is free
to fuse, and the host was fusing it while the device was not. A claim of
identical numerics that holds only when the compiler happens not to contract is
a much weaker claim than it appears, so contraction is off on both sides. On
kernels that are bandwidth bound this costs nothing measurable.

\section{The solver zoo, and why it is a family}

Nine methods are required by the objectives and three more earn their place.
Eleven of the twelve are one family: every classical stationary method is a
choice of matrix splitting $A = M - N$, and the differences between Richardson,
Jacobi, the Gauss Seidel variants, SOR and the block methods are entirely
differences in the choice of $M$. Chapter \ref{ch:background} presents them that
way, derives each in a few lines from the common framework, and cites the source
for each convergence result.

Conjugate gradient is the twelfth and is deliberately not in that list. It is
not stationary, its optimality argument is different in kind, and presenting it
as the contrast is what makes the family visible as a family.

Every convergence claim in this report is checked as a measurement. The Jacobi
spectral radius is compared against $\cos(\pi h)$; the ratio of Jacobi to
Gauss Seidel iteration counts is compared against the factor of two that Young's
theory predicts; the closed form optimal relaxation factor is tested as a
minimum by perturbing it in both directions; and the change of growth order from
$O(n^2)$ to $O(n)$ is demonstrated by measuring how the counts grow as the grid
doubles.

\section{A comparison the hardware does not decide}

A five point stencil sweep performs three memory operations and about six
floating point operations per unknown. It is bandwidth bound on every machine
ever built. A ratio of runtimes between a CPU and a GPU on such a kernel is
therefore, to within a small correction, a ratio of memory bandwidths wearing
the costume of a statement about parallel programming.

Chapter \ref{ch:methodology} sets out the protocol that avoids this. Iteration
counts come first because they are properties of the mathematics and transfer to
any machine. Efficiency against each device's own measured STREAM triad
bandwidth comes second, because it is dimensionless and says how well a device
is being used rather than which device is larger. Absolute time comes last,
labelled as a property of this specific pair of devices. The measured speedup is
then decomposed into the factor that is the memory system and the factor that is
the implementation, so a reader can see which is which.

Like is compared against like throughout. Natural ordering Gauss Seidel cannot
run on a GPU at all, because it is sequentially dependent; the comparison
therefore pairs GPU Jacobi against CPU Jacobi and GPU red black Gauss Seidel
against CPU red black Gauss Seidel, and reports the cost of the red black
reordering separately so that it is charged to both sides equally.

\section{What the environment allowed, and what it did not}

The target machine has an Intel Core i7-14700K with eight performance cores and
twelve efficiency cores, and everything runs inside WSL2. The guest does not
pass the heterogeneity through: it reports fourteen uniform cores, identical
cache sizes on every logical processor, and no hybrid feature flag. Thread
affinity binds to a guest virtual processor that the hypervisor may place on any
host core.

Rather than assume a mapping that cannot be verified from inside the guest, the
library measures. A probe times an identical kernel on each logical processor
and reports a two group split only when the throughputs separate by a margin
that dominates the noise. The core knee that the objectives ask for is taken
from the aggregate scaling curve, which depends only on how many cores are
engaged and not on which, and therefore survives virtualisation. Section
\ref{sec:knee} reports both what could and what could not be established.

\section{How to read this report}

Chapter \ref{ch:background} derives the solver family and states the convergence
theory. Chapter \ref{ch:parallelisation} describes the backend interface, what
each implementation does underneath, and the communication analysis for the
distributed backends. Chapter \ref{ch:methodology} is the comparison protocol,
written to be difficult to misquote. Chapter \ref{ch:results} contains the
measurements. Chapter \ref{ch:discussion} interprets them, including the parts
that did not come out as expected. Chapter \ref{ch:conclusion} summarises what
the work established and what it did not.

Every number in Chapter \ref{ch:results} is generated from
\texttt{experiments/results/summary.csv} by a script, and every row of that file
carries the backend, worker count, pinning policy, seed and commit hash that
produced it. A number that has not been measured reads \texttt{pending}.
